% ch02e —— 原书 p.87–93（§2.5 后段至章末）；p.94 为空白页
% 本块承接 ch02d：命题 2.61 的陈述已在 ch02d 写完，这里从 p.87 的证明开始。

% [原文 p.87]
\begin{proof}
若 $D$ 是 nef 且大的，则由引理 2.60，存在十分丰富除子 $A$ 与有效除子 $E$，使 $mD\sim A+E$。用关于 $n$ 的归纳法证明 $m^n(D^n)\ge(A^n)$。此时有等式
\[
\begin{aligned}
m^n(D^n)&=(A\cdot(mD)^{n-1})+m^{n-1}(E\cdot D^{n-1})\\
&=m^{n-1}\bigl[((D|_A)^{n-1})+((D|_E)^{n-1})\bigr].
\end{aligned}
\]
这里由定理 1.38，$((D|_E)^{n-1})$ 非负。由引理 2.60，$D|_A$ 是 nef 且大的。由归纳法，第 1 项满足下述不等式：
\[
m^{n-1}((D|_A)^{n-1})\ge((A|_A)^{n-1})=(A^n).
\]
由此可知，由 (1) 可推出 (2)。

由 (1) 推出 (3) 可由下述等式看出：
\[
D=\frac1k\bigl(A+(k-m)D\bigr)+\frac1kE.
\]
这里注意，由定理 1.18，对 $k\ge m$，$A+(k-m)D$ 是丰富的。

下面假定 $D$ 是 nef 的，且满足 $(D^n)>0$。假设基域为 $\C$，证明 $D$ 是大的。一般情形需要另外的论证，详见 [Fuj83] 与 [Kol96, VI.2.15]。

为证明大性，把所有对象经双有理态射 $f:X'\to X$ 拉回后再证明即已足够。因此可设 $X$ 非奇异。设 $B\subset X$ 为有效丰富除子，且 $B-K_X$ 也丰富。由 Riemann–Roch 弱定理 1.36 与定理 2.47，可得
\[
h^0(X,\cO(mD+B))=\chi(X,\cO(mD+B))=\frac{(D^n)}{n!}m^n+O(m^{n-1}).
\]
由正合列
\[
0\to H^0(X,\cO(mD))\to H^0(X,\cO(mD+B))\to H^0(B,\cO(mD+B|_B))
\]
得到下式：
\[
h^0(X,\cO(mD))\ge h^0(X,\cO(mD+B))-h^0(B,\cO(mD+B|_B))
=\frac{(D^n)}{n!}m^n+O(m^{n-1}).
\]
% [原文 p.88]
最后假定 (3) 成立。此时 $D=\lim A_k$ 是 nef 的，故由引理 2.60 也是大的。同样，由 (4) 可推出 (1)。

为证明由最初 3 个性质可推出最后一个性质，首先任取奇点消解 $g:X'\to X$，使 $g^{-1}(\Delta)$ 成为除子。注意：若 (4) 对 $X'$、$g^*D$ 与 $g^{-1}(\Delta)$ 成立，则 (4) 对 $X$、$D$ 与 $\Delta$ 也成立。因此可设 $X$ 自身非奇异。

由引理 2.60，可用丰富 $\Q$-除子 $A$ 与有效 $\Q$-除子 $E$ 表示成 $D\equiv A+E$。设 $f:Y\to X$ 为 $(X,E+\Delta)$ 的对数奇点消解。此时 $f^*D\equiv f^*A+f^*E$，且 $f^*E$ 是 snc 除子，但要注意 $f^*A$ 不是丰富的。由引理 2.62，存在有效 $f$-例外除子 $F$，使 $-F$ 为 $f$-丰富。此时 $A':=f^*A-\varepsilon F$ 对 $0<\varepsilon\ll 1$ 是丰富的（命题 1.45），$E':=f^*E+\varepsilon F$ 是有效 snc 除子。于是
\[
A'_k=f^*D-\frac1kE'=\frac{k-1}{k}f^*D+\frac1kA'
\]
是丰富的。
\end{proof}

\begin{lemmap}{2.62}
设 $f:Y\to X$ 为双有理态射。$Y$ 是射影的，$X$ 是 $\Q$-分解的。此时存在有效 $f$-例外除子 $F$，使 $-F$ 为 $f$-丰富。
\end{lemmap}

\begin{proof}
设 $H$ 为 $Y$ 上十分丰富的有效除子。由关于 $X$ 的假设，这样的 $H$ 是存在的。$f_*H$ 是 $\Q$-Cartier 除子，故 $f^*(f_*H)$ 有定义。对适当的有效 $f$-例外除子 $F$，成立 $f^*(f_*H)=H+F$；此时 $-F\equiv_f H$。
\end{proof}

\begin{corollaryp}{2.63}
记号与假设同引理 2.62。此时 $f$ 的例外集 $\Exc(f)$ 具有纯余维数 1。
\end{corollaryp}

\begin{proof}
设 $H,F$ 与引理 2.62 的证明中相同。$H=f^*(f_*H)-F$ 在 $Y$ 上十分丰富。由于 $f^*f_*\cO_Y(-F)\to\cO_Y(-F)$ 是满射，作为集合有 $F=f^{-1}f(F)$。因此在 $Y\setminus\Supp F=f^{-1}(X\setminus f(\Supp F))$ 上，$f^*(f_*H)=H$ 是十分丰富的，且 $f$ 是同型态射。故 $\Exc(f)=\Supp F$。
% [原文 p.89]
[Sha94, II.4.4] 中的讨论可以推广，使之也适用于这一设定。
\end{proof}

至此，为陈述并证明前面约定的定理 2.47 的推广，准备工作已经完成。

\begin{theoremp}[小平消没定理 II，{[Kaw82, Vie82]}]{2.64}
设 $X$ 为非奇异射影簇。$L$ 为 $X$ 上的线丛，且满足 $L\equiv M+\sum a_iD_i$。其中
\begin{enumerate}[label=（\arabic*）,leftmargin=2.8em]
\item $M$ 是 nef 且大的 $\Q$-除子，
\item $\sum D_i$ 是 snc 除子，
\item 对所有 $i$ 成立 $0\le a_i<1$，且 $a_i\in\Q$。
\end{enumerate}
此时对所有 $i<\dim X$，成立 $H^i(X,L^{-1})=0$。
\end{theoremp}

% [Tutor Note: 系数条件 $0\le a_i<1$ 是 klt 性的系数版本；正因 $\sum a_iD_i$ 的系数严格小于 1，其整数部分 $\lfloor\sum a_iD_i\rfloor$ 恒为 0，定理 2.64 才能逐项归约到定理 2.47 的消没结论。]

\begin{proof}
把定理 2.64 分成若干步，归约到小平消没定理。

\textbf{步骤 1}\quad 设 $f:Y\to X$ 为 $n$ 维正规簇之间的有限态射。若特征为 0，则 $(1/\deg f)\operatorname{Trace}_{Y/X}$（参看定义 5.6）给出包含映射 $\cO_X\to f_*\cO_Y$ 的分解（splitting）。因此，若 $F$ 是 $X$ 上的凝聚层，则 $F$ 是 $f_*(f^*F)$ 的直和因子。因此 $H^i(X,F)$ 是 $H^i(Y,f^*F)$ 的直和成分。这可反复应用于 $F=L^{-1}$。

\textbf{步骤 2}\quad 本步用关于除子 $D_i$ 个数的归纳法，把定理 2.64 归约到 $L=M$ 的情形。

$D_i$ 的系数是有理数，故 $L-a_1D_1$ 不是 Cartier 除子，这一归约并不显然。用步骤 1 与命题 2.67，归约到 $a_1D_1$ 数值等价于 Cartier 除子的情形。为便于归纳法论证，证明下述在更一般设定下成立的结果。

\par\medskip\noindent\textbf{主张 2.65}\quad 设 $X$ 为非奇异复簇，$L$ 为 $X$ 上的线丛，且满足 $L\equiv M+\sum a_iD_i$。其中 $D_i$ 是非奇异的（可能可约的）除子，$\sum D_i$ 是 snc 除子，$0\le a_i<1$ 是有理数。此时存在从非奇异射影簇 $Z$ 出发的满射有限态射 $p:Z\to X$，以及 $Z$ 上的线丛 $M_Z$，使 $M_Z\equiv p^*M$，并且 $H^i(X,L^{-1})$ 是 $H^i(Z,M_Z^{-1})$ 的直和因子。

\begin{proof}
对除子 $D_i$ 的个数用归纳法证明。把有理数 $a_1$ 写成 $a_1=b/m$，其中 $m\in\Z_{>0}$。由命题 2.67，存在满射有限态射 $p_1:X_1\to X$ 与 $X_1$ 上的 Cartier 除子 $D$，使 $p_1^*(D_1)\sim mD$。还可选取使各 $p_1^*(D_i)$ 非奇异、$\sum p_1^*(D_i)$ 为 snc 除子。由步骤 1，$H^i(X,L^{-1})$ 是 $H^i(X_1,p_1^*L^{-1})$ 的直和因子。$D_1$ 对应 $\cO_X(mD)$ 的截面，设 $p_2:X_2\to X_1$ 为对应的循环覆盖。由引理 2.51，$X_2$ 非奇异；同样可证各 $p_2^*(D_i)$ 非奇异，且 $p_2^*(\sum_{i>1}D_i)$ 是 snc 除子。
% [原文 p.90]
此时 $(p_2)_*\cO_{X_2}=\bigoplus_{j=0}^{m-1}\cO_{X_1}(-jD)$ 成立，故
\[
H^i(X_2,p_2^*p_1^*L^{-1}(bD))=\bigoplus_{j=0}^{m-1}H^i(X_1,p_1^*L^{-1}((b-j)D)).
\]
由 $j=b$ 的情形可知，$H^i(X_1,p_1^*L^{-1})$ 是 $H^i(X_2,p_2^*p_1^*L^{-1}(bD))$ 的直和因子。

由 $p_2^*p_1^*L(-bD)\equiv p_2^*p_1^*M+\sum_{i>1}a_ip_2^*p_1^*(D_i)$，把 $X_2$ 取作新的 $X$、$p_2^*p_1^*L(-bD)$ 取作新的 $L$，则它们满足主张 2.65 的假设。因此由关于 $\sum a_iD_i$ 项数的归纳法得到 $Z\to X_2$。合成态射 $Z\to X_2\to X$ 满足主张 2.65 的条件。
\end{proof}

\textbf{步骤 3}\quad 在定理 2.64 中假定 $M$ 是丰富的。此时应用上一步，定理 2.64 被归约到定理 2.47。因此这一情形的证明已经完成。

\textbf{步骤 4}\quad 设 $L$ 是 nef 且大的线丛。需证明对 $i<\dim X$ 有 $H^i(X,L^{-1})=0$。当 $i>0$ 时，由 Serre 对偶，这等价于 $H^i(X,\omega_X\otimes L)=0$。

取 $X$ 上任一丰富除子 $H$。由引理 2.60(4)，存在非奇异射影簇 $Y$ 与双有理态射 $f:Y\to X$，使 $f^*L\equiv A+E$。其中 $A$ 是丰富的，$E=\sum e_jE_j$ 是 snc 除子，且对任一 $j$ 有 $0\le e_j<1$。此时有下述 Leray 谱序列：
\[
H^i(X,L(rH)\otimes R^jf_*\omega_Y)\Rightarrow H^{i+j}(Y,\omega_Y\otimes f^*L(rH)).
\]
$f^*L(rH)\equiv(A+rf^*H)+E$，且 $A+rf^*H$ 是丰富的。因此如上一阶段所证，对 $r\ge 0$、$k>0$ 成立 $H^k(Y,f^*L(rH)\otimes\omega_Y)=0$。
% [原文 p.91]
首先选取 $r\gg 1$。此时对 $i>0$ 有 $H^i(X,L(rH)\otimes R^jf_*\omega_Y)=0$，且由谱序列得
\[
H^0(X,L(rH)\otimes R^jf_*\omega_Y)=H^j(Y,\omega_Y\otimes f^*L(rH))=0.
\]
由于已选 $r\gg 1$，由此可知对 $j>0$ 成立 $R^jf_*\omega_Y=0$。

最后令 $r=0$。此时得到
\[
H^i(X,L\otimes f_*\omega_Y)=H^i(Y,f^*L\otimes\omega_Y)=0.
\]
至此定理 2.64 的证明完成。
\end{proof}

\begin{remark}{2.66}
由 Serre 对偶，$H^i(X,L^{-1})$ 的消没定理与 $H^{n-i}(X,\omega_X\otimes L)$ 的消没定理等价；但在定理 2.64 的证明途中，这两个命题表现出不同的行为。步骤 1–3 处理的是 $H^i(X,L^{-1})$ 的消没定理。在同一步中处理对偶版的消没定理是可能的，但需要仔细比较 $\omega_Z$ 与 $\omega_X$，会更困难。

最后一步处理的是 $H^{n-i}(X,\omega_X\otimes L)$ 的消没定理，这是因为在这里处理 $H^i(X,L^{-1})$ 会更困难。
\end{remark}

证明途中用了 [BG71] 的下述引理，这是一个在许多场合都有用的结果。

\begin{propositionp}{2.67}
设 $X$ 为拟射影簇，$D$ 为 $X$ 上的 Cartier 除子，$m$ 为自然数。此时存在正规簇 $Y$、满射有限态射 $g:Y\to X$，以及 $Y$ 上的 Cartier 除子 $D'$，满足 $g^*D\sim mD'$。

若 $X$ 非奇异，且 $\sum F_j$ 是 $X$ 上的 snc 除子，则还可选取使 $Y$、$g^*F_j$ 非奇异，且 $\sum g^*F_j$ 为 snc 除子。
\end{propositionp}

\begin{proof}
设 $n=\dim X$，定义态射 $\pi:\PP^n\to\PP^n$ 为 $(x_0:\dots:x_n)\mapsto(x_0^m:\dots:x_n^m)$。此时 $\pi^*\cO_{\PP^n}(1)\cong\cO_{\PP^n}(m)$。

设 $L$ 为 $X$ 上十分丰富的除子。此时存在态射 $h:X\to\PP^n$，满足 $L\cong h^*\cO_{\PP^n}(1)$。设 $Y$ 为纤维积的正规化，使下述图式交换：
% [原文 p.92]
\begin{tikzcd}
Y \arrow[r,"h_Y"] \arrow[d,"g"'] & \PP^n \arrow[d,"\pi"] \\
X \arrow[r,"h"'] & \PP^n
\end{tikzcd}

此时 $g^*\cO_X(D)\cong h_Y^*(\pi^*\cO_{\PP^n}(1))\cong h_Y^*\cO_{\PP^n}(m)$。
若 $D$ 十分丰富，则最初部分的证明至此完成。
若 $X$ 非奇异，则不用 $\pi$，而用 $\PP^n$ 的一个一般自同构 $\sigma$ 与 $\pi$ 的复合 $\pi'=\pi\circ\sigma$。由 Kleiman 的 Bertini 型定理（[Har77, III.10.8]），$Y$ 非奇异，且各 $g^*F_j$ 非奇异、$\sum g^*F_j$ 是 snc 除子。
一般情形下，把 $D$ 写成 $D\sim L_1-L_2$。把上述论证用两次，得到态射 $g:Y\to X$，使 $g^*L_i\sim mL'_i$（$L'_i$ 为适当的 Cartier 除子）。令 $D':=L'_1-L'_2$ 即可。
\end{proof}

下述结果是定理 2.64 的特别情形，但作为推论陈述仍有价值。特别地，$L=0$ 的情形最初由 [GR70] 证明。

\begin{corollaryp}{2.68}
设 $f:Y\to X$ 为射影簇之间的双有理态射，$Y$ 非奇异。$L$ 为 $Y$ 上的线丛，且 $L\equiv M+\sum a_iD_i$。其中假定下述 3 个条件：
\begin{enumerate}[label=（\arabic*）,leftmargin=2.8em]
\item $M$ 是 nef 的，
\item $\sum D_i$ 是 snc 除子，
\item 对所有 $i$，$a_i\in\Q$ 满足 $0\le a_i<1$。
\end{enumerate}
此时对 $i>0$ 成立 $R^if_*(\omega_Y\otimes L)=0$。特别地，对 $i>0$ 成立 $R^if_*\omega_Y=0$。
\end{corollaryp}

\begin{proof}
设 $H$ 为 $X$ 上的丰富除子。把定理 2.64 应用于 $Y$ 上的 $L(rf^*H)$，再由命题 2.69 即得。
\end{proof}

% [Tutor Note: 页图与英文底本此处均作 $L(rf^*H)$，$r$ 是定理 2.64 的步骤 4 中选取的大参数（那里把定理 2.64 用于 $L(rH)$，$L(rH)=L\otimes\cO_X(rH)$），此处对 $Y$ 上的线丛 $L$ 取 $L(rf^*H)$ 并令 $r\gg1$，并非笔误。推论 2.68 是相对消没定理的原型，第 3 章证明极射线收缩与锥定理时反复用到。]

\begin{propositionp}{2.69}
设 $f:Y\to X$ 为固有态射，$F$ 为 $Y$ 上的凝聚层。对任一 $j$，下述条件等价：
\begin{enumerate}[label=（\arabic*）,leftmargin=2.8em]
\item 对十分丰富的 $H$，$H^j(Y,F\otimes f^*H)=0$；
\item $R^jf_*F=0$。
\end{enumerate}
\end{propositionp}

\begin{proof}
这已由上面的步骤 4 本质地处理完毕。选取 $H$，使对所有的 $i>0$ 与 $k$ 有 $H^i(X,H\otimes R^kf_*F)=0$。此时 Leray 谱序列在 $E_2$ 项退化。因此 $H^i(Y,F\otimes f^*H)=H^0(X,H\otimes R^if_*F)$。
\end{proof}

% [原文 p.93]
定理 2.64 还可进一步推广到下述形式。

\begin{theorem*}[一般小平消没定理]{2.70}
设 $(X,\Delta)$ 为固有 klt 对。$N$ 为 $X$ 上的 $\Q$-Cartier Weil 除子，且 $N\equiv M+\Delta$。其中 $M$ 是 nef 且大的 $\Q$-Cartier $\Q$-除子。
此时对 $i<\dim X$ 有 $H^i(X,\cO_X(-N))=0$。
\end{theorem*}

把定理 2.70 归约到定理 2.64 并不困难。这是 [Kol95, Chs.9–11] 中采用的手法。用正特征的另一个证明见 [EV92]。关于进一步推广的近期解说，请参看 [Kol97, sec.2]。

最后，就曲面情形的消没定理而言，如下引理所表明，snc 条件往往是不必要的。

\begin{propositionp}{2.71}
记号同推论 2.68，并假设 $Y$ 是曲面。此时去掉 (2) 的 snc 条件，结论仍成立。
\end{propositionp}

\begin{proof}
设 $g:Y'\to Y$ 为 $(Y,\sum D_i)$ 的对数奇点消解，记 $E=g^*(\sum a_iD_i)-g_*^{-1}(\sum a_iD_i)$。此时 $E$ 是 $g$-例外的有效 $\Q$-除子，且成立 $g^*L-\lfloor E\rfloor\equiv g^*M+\sum a_ig_*^{-1}D_i+\{E\}$。由 Leray 谱序列，有自然单射
\[
R^1f_*\bigl(g_*(\omega_{Y'}(-\lfloor E\rfloor))\otimes L\bigr)\subset R^1(f\circ g)_*\bigl(\omega_{Y'}\otimes g^*L(-\lfloor E\rfloor)\bigr).
\]
把推论 2.68 应用于 $f\circ g$ 与 $g^*L-\lfloor E\rfloor$，得到
\[
R^1f_*\bigl(g_*(\omega_{Y'}(-\lfloor E\rfloor))\otimes L\bigr)=R^1(f\circ g)_*\bigl(\omega_{Y'}\otimes g^*L(-\lfloor E\rfloor)\bigr)=0.
\]
单射 $g_*(\omega_{Y'}(-\lfloor E\rfloor))\otimes L\to\omega_Y\otimes L$ 的余核（cokernel）具有 0 维支集。因此得 $R^1f_*(\omega_Y\otimes L)=0$。
\end{proof}

% [Tutor Note: 曲面情形下 $\lfloor E\rfloor$ 只剩曲线，余核的支集被压成 0 维，故可用 Leray 谱序列直接消去 snc 假设；维数再高这一步就不成立了。]

% [原文 p.94] 原书此页空白，无正文。

\chapterreview{第 2 章以“极小模型纲领能否执行”为主线，把第 1 章在曲面与三维中具体操作过的收缩、翻转提升为一般的纲领。核心问题是：给定非奇异射影簇 $X$，能否通过一系列保持双有理等价的可控操作得到标准类 $K$ 为 nef 的极小模型；若 $K$ 非 nef，则得到 Fano 纤维空间。

知识脉络上，§2.1 由例 2.1 的三维商奇点说明三维以上不能只在非奇异簇内工作，随即陈述锥定理、极射线的三种收缩与翻转，给出极小模型的定义与森纲领的四步流程。§2.2 转向若干不需要高深概念的推广，其中“簇 + 除子”的对（log 范畴）最有用。§2.3 因此系统引入对 $(X,\Delta)$ 的奇点类，即标准、终端、klt、dlt，并以差异系数 $a(E,X,\Delta)$ 统一定义，进而证明 dlt 的稳定性（命题 2.42、定理 2.44）。§2.4–2.5 处理纲领所依赖的消没定理：没有它，翻转的存在性与相对收缩都无法论证。

核心结论有四条。其一，命题 2.42 与定理 2.44 把 dlt 的几何条件翻译成差异系数的数值条件，使 dlt 在消解、收缩与翻转之下保持。其二，小平消没定理 2.47：非奇异射影簇上的丰富线丛 $L$ 满足 $H^i(X,L^{-1})=0$（$i<\dim X$）；证明走 Hodge 理论、GAGA 与循环覆盖的路线，把全纯消没归约为拓扑事实与 Serre 消没。其三，定理 2.64（Kawamata–Viehweg）：$L\equiv M+\sum a_iD_i$、$M$ nef 且大、$\sum D_i$ 为 snc 除子、$0\le a_i<1$ 时同样的消没成立；它是相对消没定理的原型，推论 2.68 是其相对形式。其四，定理 2.70 把它推广到 klt 对 $N\equiv M+\Delta$，定理 2.71 说明曲面情形下 snc 条件可以去掉。围绕这些结论，大、nef 与“接近丰富”被系统化（定义 2.59、引理 2.60、命题 2.61）。

辨析与提醒。第一，丰富、nef、大三者不可混同：nef 是对所有曲线取非负，大是截面增长达到 $k^n$ 的量级，二者合起来才等价于可写成丰富 $\Q$-除子与小的有效除子之和（命题 2.61）。第二，klt 与 dlt 的差别在系数：klt 要求所有差异系数小于 1，dlt 只要求存在简单正规交叉的对数消解并允许系数等于 1。第三，消没定理的系数条件 $0\le a_i<1$ 不能放宽到 1，否则 klt 性失效，消没结论一般不再成立。

评注。本章把“纲领能否进行”分解为奇点分类与消没定理两个技术支柱，这一分解直接决定了第 3 章的次序：先用相对消没证明极射线收缩与锥定理，再讨论翻转的存在性。历史上，Kawamata–Viehweg 消没（[Kaw82, Vie82]）与 Kollár 的循环覆盖论证（[Kol95]）把小平的古典结果推广到对数情形；三维翻转的存在性由森重文于 1988 年证明，四维以上至今仍是瓶颈。读法建议：§2.3 的定义先只掌握终端、klt、dlt 三档与差异系数的符号约定，细节可在第 3 章回头查；§2.4–2.5 的证明可先跳过 GAGA 与谱序列的技术段落，抓住“用拓扑控制全纯”这条主线。}
